%% two-spines-converge.tex — the book's conceptual climax (Chapter 10):
%% the two spines enter the evolutionary loop. House conceptual style.
%% _concept-style.tex: shared TikZ styles for the course's CONCEPTUAL block diagrams.
%% The leading underscore marks this as the one file in figures/ that is not a
%% figure; every other figures/*.tex holds exactly one figure body.
%%
%% \input this at the top of every conceptual figure (before \begin{tikzpicture})
%% so each such diagram speaks one visual language: rounded "block" nodes, stealth
%% arrows, and natural `<hue>!15' fills. Redefining these styles on each \input is
%% idempotent and harmless. The reference for the house parameters is
%% resources/STYLE-figures.md. Figure~\ref{figure:AgentLoop} is the exemplar.
%%
%% Scope: use for conceptual/relational diagrams (boxes-and-arrows). Do NOT use for
%% product mock-ups (the rig figures have their own shared language) or for
%% data/plot figures.
\tikzset{
  % the standard block node: geometry only — apply a `fill=<hue>!15' per node.
  conceptbox/.style={draw, rounded corners=4pt, minimum width=2.4cm,
                     minimum height=0.85cm, align=center, font=\small},
  % a flow arrow between blocks.
  conceptflow/.style={->, >=stealth', thick},
  % an edge/annotation label.
  conceptlabel/.style={font=\scriptsize, align=center},
}%

\begin{tikzpicture}[
    spine/.style={->, >=stealth', thick, densely dashed},
]
% --- the two spine sources (left) ---
\node[conceptbox, fill=green!20] (harness) at (0, 2.0)  {Evaluation\\Harness};
\node[conceptbox, fill=blue!15]  (context) at (0,-2.0)  {Context\\Engineering};
\node[conceptlabel, above=1pt of harness] {Spine 1 (Chapter~\ref{chapter:Evaluation})};
\node[conceptlabel, below=1pt of context] {Spine 2 (Chapters~\ref{chapter:ContextEngineering}, \ref{chapter:Retrieval})};

% --- the evolutionary loop (right): a triangle Fitness -> Select -> Mutate ---
\node[conceptbox, fill=green!20]  (fit) at (5.6, 2.0)  {Fitness};
\node[conceptbox, fill=orange!20] (sel) at (8.0, 0.0)  {Select};
\node[conceptbox, fill=blue!15]   (mut) at (5.6,-2.0)  {Mutate};

% internal loop arrows
\draw[conceptflow] (mut) -- node[conceptlabel, left]  {candidates} (fit);
\draw[conceptflow] (fit) -- node[conceptlabel, above right, pos=0.4] {scored} (sel);
\draw[conceptflow] (sel) -- node[conceptlabel, below right, pos=0.6] {survivors} (mut);

% --- the two spines feeding the loop ---
\draw[spine] (harness) -- node[conceptlabel, above] {becomes} (fit);
\draw[spine] (context) -- node[conceptlabel, below] {drives}   (mut);
\end{tikzpicture}%
